% Generated draft: requires mathematical and source review.
\documentclass[11pt]{article}
\usepackage[margin=25mm]{geometry}
\usepackage{fontspec}
\setmainfont{Latin Modern Roman}
\usepackage{amsmath,amssymb,mathtools,unicode-math}
\setmathfont{Latin Modern Math}
\usepackage{enumitem,needspace,xurl,hyperref}
\setlength{\parindent}{0pt}
\setlength{\parskip}{5pt}
\setlength{\emergencystretch}{3em}
\newcommand{\N}{\mathbb{N}}
\newcommand{\Z}{\mathbb{Z}}
\newcommand{\Q}{\mathbb{Q}}
\newcommand{\R}{\mathbb{R}}
\newcommand{\C}{\mathbb{C}}
\newcommand{\F}{\mathbb{F}}
\newcommand{\A}{\mathbb{A}}
\newcommand{\PP}{\mathbb{P}}
\newcommand{\E}{\mathbb{E}}
\newcommand{\eps}{\varepsilon}
\newcommand{\dd}{\,\mathrm{d}}
\begin{document}
\section*{UnsolvedMath ID: 700338}
\textbf{Source problem code:} catalogue entry
\subsection{Definitions and mathematical statement}
A group is locally finite if every subgroup generated by finitely many elements is finite. For a prime \(r\), an \(r\)-subgroup is a subgroup whose elements all have finite orders that are powers of \(r\), including \(1=r^0\). A \(\{p,q\}\)-group is a group in which every element has finite order of the form \(p^a q^b\), where \(a,b\) are nonnegative integers.

Let \(p\) and \(q\) be distinct primes, and let \(G\) be a locally finite group. Suppose there exist a \(p\)-subgroup \(P\leq G\) and a \(q\)-subgroup \(Q\leq G\) such that
\[
G=PQ=\{xy:x\in P,\ y\in Q\}.
\]
Neither subgroup is assumed to be normal. Must \(G\) be a \(\{p,q\}\)-group? Explicitly, must every \(g\in G\) satisfy
\[
\operatorname{ord}(g)=p^a q^b
\quad\text{for some nonnegative integers }a,b?
\]
\subsection{Short English statement}
If a locally finite group is the setwise product of a \(p\)-subgroup and a \(q\)-subgroup for distinct primes \(p,q\), must every element have order divisible only by these two primes?
\subsection{Sources}
\begin{itemize}
\item Local UnsolvedMath registry: permanent ID \(700338\), exact source code \texttt{catalogue entry}, supplied complete source record.
\item B. Hartley, Problem 5.59, \emph{The Kourovka Notebook (Unsolved Problems in Group Theory)}, originally in the 5th issue (1976). The supplied record cites the 21st issue edition (2026), \texttt{arXiv:1401.0300v48}, and identifies E. I. Khukhro and V. D. Mazurov as editors of the collection.
\item The supplied registry labels this problem \texttt{open}. This is supplied source metadata, not an independently verified assessment of its current mathematical status.
\end{itemize}
\end{document}
